\chapter{Metodologia}\label{cap_3}

Este capítulo descreve o funcionamento do \ac{FLIM}, os três componentes avaliados sobre ele, e o protocolo experimental. A descrição é detalhada de propósito: boa parte das conclusões desta dissertação decorre de propriedades específicas do procedimento de estimação dos filtros, e elas só ficam visíveis quando o procedimento é apresentado com precisão.

\section{O FLIM em detalhe}

\subsection{O que é um filtro convolucional}

Uma rede convolucional reconhece padrões por meio de filtros. Um filtro é uma pequena matriz de números, tipicamente de tamanho $3 \times 3$, que é deslizada sobre a imagem. Em cada posição, calcula-se a soma dos produtos entre os valores do filtro e os valores da imagem naquela vizinhança. O resultado é alto onde a vizinhança se parece com o padrão guardado no filtro, e baixo onde não se parece.

Formalmente, para uma imagem $I$ e um filtro $K$ de tamanho $h \times w$ aplicado na posição $(x, y)$:

\begin{equation}
(I * K)(x, y) = \sum_{i=1}^{h} \sum_{j=1}^{w} I(x + i - c_h,\; y + j - c_w) \cdot K(i, j)
\label{eq_conv}
\end{equation}

onde $c_h$ e $c_w$ indicam o centro do filtro. O conjunto de filtros de uma camada define o vocabulário visual com que a rede descreve a imagem, e uma rede com várias camadas compõe padrões simples das primeiras camadas em padrões mais complexos nas seguintes.

A diferença entre os métodos está em como esses filtros são obtidos. Redes treinadas por retropropagação os inicializam aleatoriamente e os ajustam ao longo de milhares de iterações, calculando o quanto cada valor contribuiu para o erro e corrigindo-o na direção oposta. Esse procedimento requer grandes conjuntos anotados.

\subsection{Estimação dos filtros a partir de marcadores}

O \ac{FLIM} não ajusta filtros, ele os \emph{extrai} \cite{souza2020flim}. O especialista desenha marcadores sobre poucas imagens, indicando regiões de objeto e de fundo. O método recorta as vizinhanças centradas nos pixels marcados e usa essas vizinhanças, após agrupamento, como filtros.

Seja $M$ o conjunto de pixels marcados em uma imagem. Para cada pixel $p \in M$, define-se o \emph{patch} $\phi(p)$ como o vetor obtido empilhando os valores das bandas na vizinhança $h \times w$ de $p$:

\begin{equation}
\phi(p) \in \mathbb{R}^{h \cdot w \cdot b}
\label{eq_patch}
\end{equation}

onde $b$ é o número de bandas da entrada daquela camada. Para a primeira camada, $b$ é o número de canais da imagem; para as seguintes, é o número de filtros da camada anterior.

O procedimento ocorre em dois níveis de agrupamento.

\paragraph{Primeiro nível, por marcador.} Os marcadores são separados em componentes conexas: um traço isolado do pincel é uma componente, e dois traços que se encostam formam uma só. Para cada componente $m$, os patches dos seus pixels são agrupados por $k$-médias em $n_m$ grupos, e os centróides resultantes tornam-se candidatos a filtro:

\begin{equation}
C_m = \operatorname{kmeans}\big(\{\phi(p) : p \in m\},\; n_m\big)
\label{eq_nivel1}
\end{equation}

O parâmetro $n_m$, chamado \texttt{nkernels\_per\_marker} na implementação, é definido na arquitetura. Nos experimentos deste trabalho ele vale 3 no conjunto de parasitas, 5 no BraTS e 15 no conjunto de conjuntivite.

\paragraph{Segundo nível, global.} Os candidatos de todas as componentes e de todas as imagens anotadas são reunidos em um único conjunto $C = \bigcup_m C_m$, que então é reduzido ao número de canais de saída da camada, $n_{\text{out}}$:

\begin{equation}
K = \operatorname{kmeans}(C,\; n_{\text{out}})
\label{eq_nivel2}
\end{equation}

Os centróides resultantes, após normalização, são os filtros da camada. A saída da camada serve de entrada para o mesmo procedimento na camada seguinte, de modo que a estimação avança camada a camada sem nenhum passo de retropropagação.

\subsection{Os dois regimes do segundo agrupamento}

A Equação \ref{eq_nivel2} tem um comportamento que é central para os resultados deste trabalho. A implementação só executa o agrupamento quando há mais candidatos do que vagas:

\begin{equation}
K =
\begin{cases}
C, & \text{se } |C| \leq n_{\text{out}} \\[4pt]
\operatorname{kmeans}(C, n_{\text{out}}), & \text{se } |C| > n_{\text{out}}
\end{cases}
\label{eq_regimes}
\end{equation}

Quando o número de candidatos não excede o número de vagas, os próprios candidatos tornam-se os filtros, sem agrupamento algum, e o número de canais da camada é reduzido para $|C|$. Acrescentar um marcador, nesse regime, acrescenta filtros e deixa os anteriores intactos.

Quando o número de candidatos excede as vagas, o agrupamento precisa escolher como particionar o conjunto. Acrescentar um marcador acrescenta candidatos, a partição é refeita sobre o conjunto ampliado, e os centróides se deslocam. Nenhum filtro é preservado por construção.

A razão entre candidatos e vagas difere bastante entre os conjuntos de dados usados aqui. Com cerca de cinquenta componentes de marcador distribuídas em três imagens, o conjunto de parasitas produz aproximadamente 150 candidatos para 200 vagas, operando abaixo do limite. O BraTS produz cerca de 250 candidatos para 8 vagas, e o conjunto de conjuntivite cerca de 750 para 16. Nos dois últimos casos há dezenas de candidatos disputando cada vaga, e o agrupamento é intenso.

\subsection{O decodificador adaptativo}\label{sec_decoder}

O codificador produz mapas de ativação, não uma máscara. A conversão é feita por um decodificador adaptativo \cite{soares2026flim}, que atribui a cada canal de ativação um peso em $\{-1, +1\}$ conforme uma heurística que estima se aquele canal responde ao objeto ou ao fundo, e combina os canais por uma convolução ponto a ponto:

\begin{equation}
S(x, y) = \sum_{c=1}^{n_{\text{out}}} w_c \cdot A_c(x, y), \qquad w_c \in \{-1, +1\}
\label{eq_decoder}
\end{equation}

onde $A_c$ é o mapa de ativação do canal $c$. Os pesos são recalculados para cada imagem de entrada, o que dá ao decodificador o caráter adaptativo. Nenhum parâmetro é ajustado por gradiente.

O mapa $S$ é normalizado para o intervalo $[0, 1]$ e binarizado por limiarização de Otsu, seguida de um filtro que remove componentes conexas fora de uma faixa de área esperada para o objeto.

A Figura~\ref{fig_pipeline} acompanha uma imagem por todas essas etapas, da entrada ao erro por pixel. É útil retê-la, porque cada resultado dos capítulos seguintes é uma medida tomada sobre o último painel.

\begin{figure}[htb]
\centering
\includegraphics[width=\linewidth]{flim_pipeline.png}
\caption[O \acs{FLIM} de ponta a ponta numa imagem]{O método de ponta a ponta. A imagem de entrada passa pelo codificador, que produz o mapa de saliência; a limiarização de Otsu com filtro de área produz a máscara binária; a comparação com a anotação de referência dá o $F_\beta$ e o \ac{IoU}. O último painel decompõe o erro em falso positivo e falso negativo.}
\label{fig_pipeline}
\end{figure}

Vale também fixar o que é um marcador, já que ele é a única entrada de supervisão do método. A Figura~\ref{fig_markers} contrasta o traço desenhado por um especialista com o traço sintético usado nas campanhas que precisam de mais imagens anotadas do que as 31 para as quais existe marcação real.

\begin{figure}[htb]
\centering
\includegraphics[width=\linewidth]{flim_markers.png}
\caption[Marcador real e marcador sintético]{O que o método recebe como supervisão. À esquerda o traço real de um especialista; à direita o traço sintético gerado a partir da anotação de referência. A distinção importa porque marcador real existe apenas para 31 imagens do conjunto de parasitas, e comparações entre braços de origens diferentes mediriam quem desenhou o traço em vez do efeito em estudo.}
\label{fig_markers}
\end{figure}

\section{Métricas}

\subsection{Qualidade da segmentação}

A métrica principal é o $F_\beta$, que combina precisão e revocação com peso maior para a precisão, seguindo o trabalho de referência:

\begin{equation}
F_\beta = \frac{(1 + \beta^2) \cdot P \cdot R}{\beta^2 \cdot P + R}, \qquad \beta^2 = 0{,}3
\label{eq_fbeta}
\end{equation}

onde $P$ é a precisão e $R$ a revocação, calculadas pixel a pixel sobre a máscara binarizada.

Reporta-se também a \ac{IoU}, razão entre a interseção e a união das máscaras predita e verdadeira:

\begin{equation}
\text{IoU} = \frac{|\hat{Y} \cap Y|}{|\hat{Y} \cup Y|}
\label{eq_iou}
\end{equation}

A acurácia pixel a pixel é reportada por completude, mas não discrimina neste domínio: como os objetos ocupam pequena fração da imagem, predizer tudo como fundo já produz acurácia superior a $0{,}97$.

\subsection{Esforço de interação}

Para a avaliação interativa usa-se o \ac{NoC}, número de cliques necessários para que a \ac{IoU} ultrapasse um alvo $\tau$, com um teto $N_{\max}$ de interações:

\begin{equation}
\text{NoC@}\tau = \min\{\,n \leq N_{\max} : \text{IoU}_n \geq \tau\,\},
\quad \text{ou } N_{\max} \text{ se nenhum } n \text{ satisfaz}
\label{eq_noc}
\end{equation}

Imagens que não atingem o alvo entram na média com o valor do teto. Por isso a taxa de sucesso é reportada ao lado do \ac{NoC} em todas as tabelas: um método poderia exibir contagem baixa simplesmente por desistir antes, e a contagem isolada premiaria esse comportamento.

\section{Aprendizado ativo}
\label{sec_al_metodo}

\subsection{O problema que o aprendizado ativo resolve}

Anotar uma imagem para segmentação custa caro. Quem anota precisa ser um especialista, o traço precisa acompanhar a fronteira do objeto, e o tempo gasto cresce com a quantidade de imagens. O conjunto de parasitas usado aqui tem 1220 lâminas; anotar todas é inviável na prática e desnecessário em princípio, porque muitas são parecidas entre si e o modelo não aprende nada de novo com a segunda cópia de algo que já viu.

O \ac{AL} parte dessa observação. Em vez de anotar tudo ou sortear ao acaso, ele usa o modelo já treinado para decidir o que vale a pena anotar em seguida. O procedimento é um laço:

\begin{enumerate}
    \item treina-se um modelo com o pouco que já está anotado, o conjunto $A$;
    \item aplica-se esse modelo ao que ainda não foi anotado, o conjunto $U$, chamado \emph{pool};
    \item uma função de aquisição $s(\cdot)$ atribui uma pontuação a cada candidato de $U$;
    \item o candidato de maior pontuação vai para o especialista, que o anota;
    \item ele passa de $U$ para $A$, e o laço recomeça.
\end{enumerate}

A função de aquisição é o que distingue um método de \ac{AL} de outro. Todo o resto do laço é igual.

\subsection{O custo de pontuar o pool, e por que ele importa aqui}
\label{sec_custo_pool}

Há um detalhe que a literatura de \ac{AL} costuma tratar como secundário e que, no caso do \ac{FLIM}, é central: \textbf{pontuar o pool também custa}. A pergunta não é só quantas imagens o especialista anota, é também o que a máquina precisa fazer, e precisa saber, para decidir quais são.

Os critérios avaliados nesta dissertação se separam em três faixas de custo, e a distinção importa porque o \ac{FLIM} existe justamente para funcionar quando não se tem anotação em massa:

\begin{description}
    \item[Uma passagem pelo pool, sem anotação.] Entropia, \ac{LC} e margem precisam apenas da predição do modelo atual sobre cada imagem de $U$. O custo é uma inferência por imagem, e nenhuma anotação é consultada. É a faixa mais barata.

    \item[Várias passagens, ou atributos de todo o pool, sem anotação.] O \ac{BALD} exige um comitê, logo $N$ inferências por imagem. O CoreSet e o \ac{BADGE} exigem extrair o vetor de atributos de \emph{cada} imagem do pool e depois resolver um problema de cobertura sobre todos eles. Continuam sem consultar anotação, mas o custo cresce com o tamanho do pool, e não apenas com o orçamento.

    \item[Anotação de todo o pool.] É o caso do passo 7 do Algoritmo 1 do trabalho de referência, que escolhe a próxima imagem entre as de pior $F_\beta$. Para saber o $F_\beta$ de uma imagem é preciso ter o gabarito dela. Escolher 3 imagens entre 848 exige, portanto, as 848 anotadas.
\end{description}

Esta última linha é a razão de ser desta dissertação. O próprio artigo declara que a abordagem de seleção é supervisionada, e um procedimento que precisa do gabarito de todo o pool para escolher o que anotar derrota o propósito de um método criado para operar com três a cinco imagens anotadas. A pergunta que organiza o trabalho é, então, quanto do benefício daquele procedimento se recupera com critérios da primeira e da segunda faixa, que não consultam anotação nenhuma.

Por isso o Algoritmo 1 aparece aqui sob o nome de \textbf{baseline supervisionado com informação privilegiada}, e não como método concorrente: ele tem acesso a uma informação que os demais não têm.

Convém ser explícito quanto ao que esse braço \emph{não} é. Consultar a anotação de referência não transforma automaticamente uma heurística em limite superior de desempenho. Um limite superior exigiria que o procedimento maximizasse o desfecho avaliado, e o Algoritmo 1 maximiza outra coisa: ele escolhe a imagem de \emph{pior} $F_\beta$, que não é necessariamente a mais informativa. A medição relatada no Capítulo~\ref{cap_4} mostra que ele pode, inclusive, perder para o sorteio aleatório. Os termos \emph{teto} e \emph{limite superior} são, portanto, evitados para esse braço ao longo de todo o trabalho.

\subsection{Critérios de incerteza}
\label{sec_crit_incerteza}

A família de incerteza parte da ideia de que o modelo aprende mais com aquilo sobre o que está em dúvida. Se a predição já é confiante, anotar aquela imagem apenas confirma o que o modelo já sabe.

\paragraph{Uma ressalva de nomenclatura, antes das fórmulas.} Seja $S(x,y) \in [0,1]$ o valor que o decodificador atribui ao pixel $(x,y)$. Esse valor é o \textbf{escore de saliência normalizado}: ele resulta da normalização linear do mapa de ativação da Equação~\ref{eq_decoder} para o intervalo unitário, e é armazenado quantizado em oito bits. Ele \textbf{não é uma probabilidade calibrada}. Não provém de uma função logística ajustada, não passou por calibração, e nada garante que o valor $0{,}6$ corresponda a 60\% de chance de o pixel ser objeto.

As funções de aquisição abaixo são, portanto, aplicadas a $S$ \emph{por analogia} com a formulação probabilística, e o termo \emph{pseudo-probabilidade} é usado ao longo do texto para marcar essa distinção. Duas consequências decorrem disso e valem ser enunciadas. A primeira é que a entropia calculada sobre $S$ não deve ser lida como incerteza epistêmica calibrada: o máximo em $S = 0{,}5$ é uma propriedade de onde a normalização situa o ponto médio, e não de uma dúvida probabilística do modelo. A segunda é que as comparações entre braços permanecem válidas, porque todos os braços recebem exatamente a mesma representação, e a diferença entre eles isola o critério e não a escala.

A \textbf{entropia} mede a indecisão ponto a ponto, atingindo o máximo quando o escore é exatamente $0{,}5$:

\begin{equation}
H(x,y) = -\big[S \log S + (1-S)\log(1-S)\big]
\label{eq_entropia}
\end{equation}

O \textbf{least confidence} mede a distância à fronteira de decisão. A intuição é a mesma da entropia, com escala diferente:

\begin{equation}
\text{LC}(x,y) = 1 - |2S - 1|
\label{eq_lc}
\end{equation}

A \textbf{margem}, em classificação multiclasse, é a diferença entre as duas classes mais prováveis. No caso binário só existem duas classes, e a definição colapsa exatamente na de \emph{least confidence}. Os dois critérios foram mantidos separados nos experimentos porque a literatura os trata como distintos, mas a equivalência algébrica no caso binário precisa ser dita: \textbf{qualquer diferença observada entre eles nos resultados é ruído de execução, não efeito de critério}.

Os três critérios acima medem a incerteza de \emph{um} modelo, e essa incerteza mistura duas coisas diferentes. Um pixel pode ser ambíguo porque a imagem é genuinamente ambígua naquele ponto, ou porque o modelo está mal treinado. Anotar o segundo caso ajuda; anotar o primeiro não resolve, porque nenhuma quantidade de dados elimina ambiguidade intrínseca.

O quarto critério desta família, aqui denominado \textbf{critério inspirado em \ac{BALD} por discordância de comitê}, separa as duas usando o desacordo entre modelos. Dado um comitê de $N$ modelos com predições $p_1, \dots, p_N$:

\begin{equation}
\text{BALD} = H\!\left(\frac{1}{N}\sum_i p_i\right) - \frac{1}{N}\sum_i H(p_i)
\label{eq_bald}
\end{equation}

O primeiro termo é a incerteza total e o segundo é a parcela que todos os modelos compartilham. A diferença isola o desacordo: um ponto em que todos os modelos estão inseguros da mesma forma recebe pontuação baixa, e um ponto em que eles discordam entre si recebe pontuação alta.

A formulação original do \ac{BALD} é bayesiana: o comitê aproxima a posterior sobre os parâmetros, obtida por \emph{dropout} de Monte Carlo com várias passagens estocásticas pelo mesmo modelo, e a quantidade medida é o ganho de informação esperado sobre essa posterior. \textbf{O \ac{FLIM} não tem \emph{dropout} e sua inferência é determinística}, de modo que essa via não existe aqui, e não há posterior a aproximar.

O que este trabalho mede é, portanto, \emph{discordância entre codificadores distintos}, usando como comitê os mapas de saliência de cada um. A quantidade da Equação~\ref{eq_bald} continua bem definida e continua separando desacordo de incerteza compartilhada, mas ela \textbf{não é uma aproximação bayesiana}, e o texto evita chamá-la de \ac{BALD} sem qualificação. A leitura dos resultados correspondentes deve ser feita nesses termos.

\subsection{Critérios de diversidade}
\label{sec_crit_diversidade}

A família de diversidade parte de uma objeção à anterior. Escolher sempre o mais incerto tende a escolher repetidamente o mesmo tipo de caso difícil, e um lote de imagens parecidas entre si ensina pouco mais do que uma delas. Os critérios de diversidade ignoram a incerteza e olham a \emph{cobertura}: o objetivo é que o conjunto anotado represente bem o conjunto todo.

Cada candidato é descrito por um vetor $f(\cdot)$ obtido pela média das ativações do codificador dentro da sua máscara.

O \textbf{CoreSet} formula a seleção como cobertura geométrica, escolhendo o candidato mais distante de tudo o que já foi anotado, que é o passo guloso do $k$-center:

\begin{equation}
r^{\ast} = \arg\max_{r \notin A} \; \min_{a \in A} \; \lVert f(r) - f(a) \rVert_2
\label{eq_coreset}
\end{equation}

onde $A$ é o conjunto já anotado e $f(\cdot)$ o vetor de atributos. Lê-se de dentro para fora: para cada candidato, encontre o item já anotado mais próximo dele; depois escolha o candidato cuja menor distância é a maior de todas. Em palavras, o candidato que está na região mais desassistida do espaço de atributos.

O \textbf{medoide} faz o oposto, escolhendo o candidato mais típico, o de menor distância média a todos os demais:

\begin{equation}
r^{\ast} = \arg\min_{r} \; \frac{1}{|R|}\sum_{q \in R} \lVert f(r) - f(q) \rVert_2
\label{eq_medoide}
\end{equation}

Ele entrou nos experimentos como contraste deliberado ao CoreSet. Se a cobertura de regiões extremas fosse o que importa, o medoide deveria ir pior que o CoreSet; se o que importasse fosse representatividade, o contrário. Ter os dois permite distinguir as hipóteses, em vez de apenas constatar que um critério não funcionou.

O terceiro critério de diversidade avaliado é um \textbf{híbrido inspirado no \ac{BADGE}}, e a qualificação não é cosmética: a implementação usada aqui diverge da formulação original de uma forma que precisa ser declarada.

A ideia do \ac{BADGE} é construir um vetor que carregue incerteza e diversidade ao mesmo tempo, aplicando o $k$-means++ sobre \emph{embeddings de gradiente}. No método original, o embedding é o gradiente da perda em relação à última camada, avaliado no pseudo-rótulo $\hat{y} = \arg\max_k P_k$:

\begin{equation}
g_i^{\text{original}} = \big(P_i - e_{\hat{y}}\big) \otimes f(r_i)
\label{eq_badge_original}
\end{equation}

cuja magnitude, no caso binário, é proporcional a $|P_i - \operatorname{round}(P_i)|$, isto é, \emph{máxima} quando $P_i$ se aproxima de $0{,}5$ e próxima de zero quando o modelo está confiante. A implementação avaliada neste trabalho usa, em seu lugar:

\begin{equation}
g_i = \big(S_i - \tfrac{1}{2}\big)\, f(r_i)
\label{eq_badge}
\end{equation}

cuja magnitude é $|S_i - 0{,}5|$, ou seja, \textbf{nula} na incerteza máxima e \emph{máxima} na confiança máxima. Como o $k$-means++ sorteia pontos com probabilidade proporcional ao quadrado da distância, o efeito é favorecer os candidatos mais \emph{confiantes}, invertendo a intenção do método original.

A divergência foi identificada na auditoria do código e está registrada aqui em vez de corrigida porque este critério \textbf{não participa de nenhum resultado central} desta dissertação: ele aparece somente na tabela que reúne todas as técnicas avaliadas, e não nas campanhas de ganho marginal, de custo do aprendizado ativo ou de troca do gerador. A leitura correta da linha correspondente é, portanto, que \emph{um híbrido de diversidade ponderado por confiança} não superou o sorteio, e não que o \ac{BADGE} tenha sido avaliado. Reimplementar a Equação~\ref{eq_badge_original} e reexecutar essa comparação consta entre os trabalhos futuros do Capítulo~\ref{cap_5}.

\subsection{Por que estes critérios, e não outros}

A escolha não foi por conveniência de implementação. O conjunto avaliado cobre as posições que a literatura ocupa, com um representante de cada e com os controles que permitem distingui-las:

\begin{itemize}
    \item \textbf{incerteza pura}: entropia, \ac{LC} e margem;
    \item \textbf{incerteza corrigida por desacordo}: \ac{BALD};
    \item \textbf{diversidade pura}: CoreSet e medoide, em direções opostas;
    \item \textbf{híbrido}: \ac{BADGE};
    \item \textbf{baseline privilegiado}: o procedimento do trabalho de referência, que lê o gabarito e reproduz o passo 7 do Algoritmo 1;
    \item \textbf{piso}: sorteio aleatório, sem critério nenhum.
\end{itemize}

O piso e o baseline privilegiado são o que torna o conjunto interpretável. Sem o sorteio não se sabe se um critério ajudou ou se qualquer escolha teria ajudado. Sem o baseline privilegiado não se sabe se a dificuldade está em o critério ser imperfeito ou em a seleção de imagem não ser a alavanca certa, e é essa distinção que os resultados do Capítulo~\ref{cap_4} acabam decidindo.

Ficaram de fora os métodos que exigem retropropagação para funcionar, como os baseados em gradiente da perda ou em redes adversárias de seleção. Não é omissão: o \ac{FLIM} não tem perda nem otimizador, e esses métodos não são aplicáveis sem descaracterizar o modelo que está sendo estudado.

\subsection{Seleção de imagem e seleção de região}
\label{sec_imagem_vs_regiao}

Os critérios acima foram aplicados em dois níveis, e os dois níveis respondem a perguntas diferentes. Confundi-los é o erro mais fácil de cometer ao ler os resultados.

Na \textbf{seleção de imagem}, o candidato é a imagem inteira e a pontuação é a média do critério sobre todos os seus pixels. O que se decide é \emph{quais imagens} o especialista vai anotar. É a forma usual na literatura e é a que o Algoritmo 1 do trabalho de referência usa. A Figura~\ref{fig_al_imagem} mostra o que o critério enxerga nesse nível.

\begin{figure}[htb]
\centering
\includegraphics[width=\linewidth]{al_selecao_imagem.png}
\caption[O que o critério de incerteza enxerga no nível de imagem]{Seleção no nível de imagem. À esquerda, o mapa de entropia produzido pelo codificador; à direita, as imagens do pool com maior e menor pontuação. A pontuação de uma imagem é a média do critério sobre todos os seus pixels, de modo que o fundo, que ocupa a maior parte da área, domina o valor.}
\label{fig_al_imagem}
\end{figure}

Na \textbf{seleção de região}, a imagem é particionada em regiões e o candidato é uma região. A pontuação é a média do critério dentro dela:

\begin{equation}
s(r) = \frac{1}{|r|} \sum_{(x,y) \in r} H(x,y)
\label{eq_score_regiao}
\end{equation}

O que se decide aqui é \emph{onde}, dentro de imagens já escolhidas, o traço vai cair. O braço de região usa exatamente as mesmas imagens que o sorteio usaria e gasta exatamente a mesma quantidade de pixels marcados. A única variável é a posição. A Figura~\ref{fig_al_regiao} mostra o procedimento completo.

\begin{figure}[htb]
\centering
\includegraphics[width=\linewidth]{al_regiao_slic.png}
\caption[Aprendizado ativo no nível de região]{Seleção no nível de região, passo a passo: a imagem é particionada pelo \ac{SLIC}, cada superpixel recebe a média do critério dentro dele, as regiões mais bem pontuadas são apresentadas ao especialista, e a resposta dele vira marcadores de objeto e de fundo. As mesmas imagens e o mesmo orçamento de pixels do braço de sorteio; o que muda é a posição.}
\label{fig_al_regiao}
\end{figure}

A consequência prática é que \textbf{as duas diferenças não são comparáveis entre si}. O $\Delta$ da seleção de imagem mede escolha de imagem; o $\Delta$ da seleção de região mede posição de anotação. Somá-los ou apresentá-los na mesma coluna mediria duas coisas ao mesmo tempo.

A extensão do CoreSet e do medoide ao nível de região é contribuição deste trabalho. Na literatura esses dois critérios são de imagem por construção, pois comparam imagens inteiras. A adaptação substitui a média global das ativações pela média dentro da máscara de cada região, mantendo o mesmo espaço de atributos e o mesmo passo guloso.

\subsection{Geração dos candidatos}
\label{sec_gerador}

A pontuação só ordena a lista de candidatos que recebe, de modo que o procedimento que gera essa lista é parte do método, e não detalhe de implementação. Dois geradores são comparados nesta dissertação.

O primeiro, usual na literatura, particiona a imagem em superpixels pelo algoritmo \ac{SLIC}, e cada superpixel é um candidato. O segundo, proposto aqui, define cada candidato como um disco de raio $\rho$ centrado em um ponto do contorno predito:

\begin{equation}
\partial \hat{Y} = \hat{Y} \setminus \operatorname{erosão}(\hat{Y}), \qquad
r_i = \{(x,y) : \lVert (x,y) - c_i \rVert_2 \leq \rho,\; c_i \in \partial \hat{Y}\}
\label{eq_contorno}
\end{equation}

onde $\hat{Y}$ é a máscara predita. Os centros $c_i$ são amostrados sobre o contorno com espaçamento mínimo de $2\rho$, para que os discos se toquem sem se sobrepor. O contorno usado é o \emph{predito}, e não o verdadeiro, de modo que o procedimento não depende de anotação prévia.

A Figura~\ref{fig_gerador} mostra o contraste entre os dois geradores na mesma imagem. A diferença não é de escore: é de qual lista o escore recebe.

\begin{figure}[htb]
\centering
\includegraphics[width=\linewidth]{gerador_candidatos.png}
\caption[Superpixels e discos de contorno como geradores de candidatos]{Os dois geradores de candidatos na mesma imagem. As arestas do \ac{SLIC} seguem as bordas presentes na imagem por construção, de modo que seus pedaços tendem a ficar de um lado só da fronteira do objeto; os discos centrados no contorno cobrem os dois lados por construção. O painel da direita usa predição ideal e ilustra a geometria do gerador: com a predição real do codificador os discos caem no interior, porque o modelo sub-segmenta e o contorno predito fica por dentro do verdadeiro, e o ganho medido no Capítulo~\ref{cap_4} acontece apesar disso.}
\label{fig_gerador}
\end{figure}

Uma limitação do gerador proposto precisa ser registrada junto com ele, porque é consequência direta do seu desenho. Ele depende de haver um contorno predito: com máscara predita vazia ou cheia, o conjunto de candidatos sai vazio, enquanto o \ac{SLIC} funciona sempre. E a dependência não aparece apenas no caso extremo. Quando o contorno predito é curto, o espaçamento mínimo de $2\rho$ limita quantos centros cabem, e o gerador entrega \emph{menos} candidatos do que os pedidos. Isso foi observado em três sementes das campanhas e tem consequência estatística, tratada no Capítulo~\ref{cap_4}.

\section{Aprendizado contínuo}

\subsection{O esquecimento no FLIM}\label{sec_cl_flim}

Os métodos consolidados de \ac{CL} supõem um modelo treinado por gradiente, e atuam sobre o deslocamento dos pesos durante o ajuste. Como o \ac{FLIM} não possui otimizador, nenhum deles se aplica diretamente.

O fenômeno análogo, contudo, está na Equação \ref{eq_nivel2}. Sejam $C_t$ o conjunto de candidatos na rodada $t$ e $K_t$ o banco de filtros resultante. Acrescentar marcadores produz $C_{t+1} \supset C_t$, e como o agrupamento é refeito sobre o conjunto ampliado, não há garantia de que $K_{t+1}$ preserve qualquer elemento de $K_t$. Além disso, $K_t$ é compartilhado por todas as imagens, inclusive as não anotadas, de modo que a perturbação é global.

\subsection{O mecanismo de retenção}

A intervenção proposta interpola o banco em vez de substituí-lo. Para cada centróide novo $k \in K_{t+1}$, localiza-se o centróide anterior mais próximo e produz-se uma combinação convexa:

\begin{equation}
\tilde{k} = (1 - \alpha)\,\operatorname{anc}(k) + \alpha\,k,
\qquad
\operatorname{anc}(k) = \arg\min_{k' \in K_t} \lVert k - k' \rVert_2
\label{eq_retencao}
\end{equation}

O parâmetro $\alpha \in [0,1]$ é a taxa de plasticidade. Com $\alpha = 1$ o procedimento reproduz o \ac{FLIM} original. Com $\alpha = 0$ o banco passa a conter apenas direções já presentes em $K_t$, e nenhuma informação nova entra. Valores intermediários deslocam o banco parcialmente na direção da nova estimativa.

A âncora é tomada por centróide \emph{novo}, e não por centróide anterior. A direção importa: ancorar no sentido oposto obrigaria a saída a ter o tamanho do banco anterior, o que conflita com o ajuste automático do número de canais descrito na Equação \ref{eq_regimes} e invalida as estatísticas de normalização recalculadas a cada rodada.

\section{Aprendizado interativo}

\subsection{O usuário simulado}

Avaliar refinamento interativo exige um procedimento reprodutível que decida onde o especialista clicaria. Adota-se o protocolo consolidado por \cite{xu2016deep_arxiv}, no qual o clique vai ao centro do maior erro.

O erro é separado em dois tipos. O falso negativo, $\text{FN} = Y \setminus \hat{Y}$, indica região de objeto que o modelo não detectou e gera clique de objeto. O falso positivo, $\text{FP} = \hat{Y} \setminus Y$, indica região de fundo marcada como objeto e gera clique de fundo. Entre as componentes conexas desses dois conjuntos, escolhe-se a de maior área.

Dentro da componente escolhida, o clique \textbf{não} vai ao centróide. Vai ao ponto de máxima transformada de distância:

\begin{equation}
c^{\ast} = \arg\max_{(x,y) \in E} \; D_E(x,y)
\label{eq_clique}
\end{equation}

onde $E$ é a componente de erro e $D_E(x,y)$ é a distância euclidiana de $(x,y)$ à borda de $E$. Essa escolha não é arbitrária. O centróide de uma região côncava pode cair fora dela, o que produziria um clique com rótulo errado; e mesmo em região convexa o centróide não é necessariamente o ponto mais afastado da borda. O ponto de máxima distância está, por construção, no interior e o mais longe possível da franja, onde a atribuição de rótulo seria ambígua.

O clique é materializado como um disco de raio 5 pixels, valor medido nos 31 conjuntos de marcadores reais disponíveis, recortado de modo a não ultrapassar a fronteira do objeto. Não há sorteio em nenhuma etapa do laço, de forma que o \ac{NoC} é determinístico dado o codificador inicial.

A Figura~\ref{fig_clique} percorre esse procedimento numa imagem real e mostra por que a distinção entre centróide e máxima distância não é preciosismo. No exemplo, a região de erro tem forma de anel, e o centróide cai no vazio central, isto é, \emph{fora} da região que deveria rotular.

\begin{figure}[htb]
\centering
\includegraphics[width=\linewidth]{il_clique.png}
\caption[Posicionamento do clique do usuário simulado]{Como o clique é posicionado. A diferença entre a predição e a anotação de referência define a região de erro; dentro da maior componente conexa dela, o clique vai ao máximo da transformada de distância. Neste exemplo a região de erro é um anel e o centróide cai fora dela, que é exatamente o caso que a Equação~\ref{eq_clique} evita.}
\label{fig_clique}
\end{figure}

\subsection{O laço}

A cada iteração, o modelo prediz, o usuário simulado escolhe o clique, o marcador é acrescentado, o codificador é reestimado com a plasticidade configurada, e a \ac{IoU} é medida novamente. O laço termina quando a \ac{IoU} ultrapassa o alvo ou quando o teto de cliques é atingido.

O laço possui dois pontos em que os componentes se encaixam. O \ac{AL} pode reordenar os candidatos a clique, em lugar de sempre tomar o maior erro, e o \ac{CL} controla como o clique é incorporado, pelo parâmetro $\alpha$ da Equação \ref{eq_retencao}.

\section{A plataforma interativa}

Os experimentos descritos neste capítulo são executados em lote, com usuário simulado. Para observar o comportamento do sistema com uma pessoa real no laço, foi desenvolvida uma aplicação que roda localmente e integra os três componentes em um fluxo único.

A aplicação é um servidor em Python servindo uma página única, sem dependência de serviço externo. Ela é iniciada em linha de comando e atendida em \texttt{localhost}, de modo que as imagens, que são dados de pacientes em dois dos três conjuntos, não saem da máquina. A interface se organiza em cinco abas, descritas a seguir.

\subsection{Os conjuntos disponíveis e o envio de novos}

A aba de conjuntos lista o que está disponível e permite enviar um conjunto novo, conforme a Figura~\ref{fig_plat_datasets}. O padrão adotado é um diretório com \texttt{orig/} e, opcionalmente, \texttt{label/}. A anotação de referência é opcional de propósito: um conjunto trazido por um especialista normalmente não tem anotação nenhuma, que é precisamente a situação que o método endereça. Sem \texttt{label/} a aplicação continua treinando e prevendo, e o que deixa de existir é o cálculo automático do $F_\beta$, passando a avaliação a ser o julgamento visual de quem usa.

\begin{figure}[htb]
\centering
\includegraphics[width=\linewidth]{plataforma_datasets.jpg}
\caption[Aba de conjuntos de dados da plataforma]{A aba de conjuntos. Além dos três conjuntos usados nos experimentos, a aplicação aceita o envio de um conjunto novo, com ou sem anotação de referência, e permite separar o lote entre treino e validação no momento do envio.}
\label{fig_plat_datasets}
\end{figure}

\subsection{O ciclo}

A aplicação opera em rodadas. A cada rodada, o estado consiste no conjunto de marcadores acumulados, no codificador estimado a partir deles, e na curva de qualidade registrada até ali.

O \ac{AL} atua na etapa de sugestão: a aplicação calcula o mapa de probabilidade, pontua as regiões candidatas, e apresenta ao especialista a de maior incerteza. O especialista decide se aceita a sugestão ou se marca outro ponto, e informa se a região é objeto ou fundo. O \ac{IL} atua na incorporação: o toque é materializado como disco de raio 5, acrescentado ao conjunto de marcadores, e o codificador é reestimado. O \ac{CL} atua no controle da reestimação, pelo parâmetro de plasticidade da Equação \ref{eq_retencao}, que pode ser alterado \emph{entre} rodadas.

Essa última possibilidade é deliberada. Alterar a plasticidade no meio de uma sessão permite comparar os dois comportamentos sobre a mesma imagem e o mesmo histórico de cliques, o que torna a diferença visível sem exigir duas sessões separadas.

A Figura~\ref{fig_plat_anotar} mostra a tela desse ciclo. À esquerda ficam o pincel de objeto, o pincel de fundo e a imagem que o critério escolheu; à direita, o estado da rodada, o orçamento configurado, o decodificador e o critério em uso, e um registro textual que explica, a cada passo, \emph{por que} aquela imagem foi escolhida. Esse registro é deliberado: sem ele a escolha do sistema é uma caixa preta para quem anota, e a aplicação existe para tornar o mecanismo observável.

\begin{figure}[htb]
\centering
\includegraphics[width=\linewidth]{plataforma_anotar.jpg}
\caption[Tela de anotação e treino da plataforma]{A tela de anotação. O painel da direita mostra a rodada corrente, o motivo da escolha da imagem, o orçamento, o decodificador e o critério de \ac{AL}. Abaixo dele fica o painel de regiões confusas, pelo qual o modelo pergunta e o especialista decide.}
\label{fig_plat_anotar}
\end{figure}

\subsection{A comparação entre braços na própria tela}

Duas comparações do Capítulo~\ref{cap_4} estão implementadas na aplicação, e é por elas que ela deixa de ser apenas uma demonstração de interface.

A primeira, na Figura~\ref{fig_plat_comparar}, põe o \ac{FLIM} com sorteio de imagens contra um ou mais critérios de \ac{AL}, com o mesmo orçamento e o mesmo conjunto de validação. Um detalhe do desenho merece nota: a imagem da primeira rodada é comum a todos os braços, e o traço feito nela vale para todos. Sem isso, a diferença observada no fim incluiria qual imagem calhou de sair primeiro para cada braço, que é a maior fonte de variância do método.

\begin{figure}[htb]
\centering
\includegraphics[width=\linewidth]{plataforma_comparar.jpg}
\caption[Comparação entre \acs{FLIM} puro e Active Learning na plataforma]{A comparação entre braços. Todos recebem o mesmo orçamento de imagens e são medidos no mesmo conjunto de validação. Quando dois critérios escolhem a mesma imagem, ela aparece uma vez só na fila e a anotação vale para ambos.}
\label{fig_plat_comparar}
\end{figure}

A segunda, na Figura~\ref{fig_plat_onde}, é a que corresponde à contribuição central desta dissertação. Ela mantém as mesmas imagens em todos os braços e o mesmo número de pixels marcados, e varia apenas \emph{quem escolhe o lugar}: o especialista, o modelo por pré-marcação, o critério de \ac{AL} por região, ou o sorteio. É o braço de sorteio que torna os outros interpretáveis, porque sem ele não há como separar o ganho de escolher bem a região do ganho de simplesmente anotar em blocos.

\begin{figure}[htb]
\centering
\includegraphics[width=0.62\linewidth]{plataforma_onde_anotar.jpg}
\caption[Comparação de onde anotar, na plataforma]{A comparação de posição de anotação. As mesmas imagens, o mesmo orçamento de pixels, e quatro formas de decidir onde o traço cai. O braço de região sorteada é o controle.}
\label{fig_plat_onde}
\end{figure}

\subsection{O que a tela registra}

Além da máscara corrente, a aplicação mantém a curva de qualidade por clique e três indicadores: a qualidade inicial, a melhor qualidade alcançada, e o número de cliques que \emph{pioraram} o resultado. O último indicador é o motivo de a tela existir. Com plasticidade máxima, que corresponde ao comportamento original, ele cresce ao longo da sessão, e o especialista observa diretamente o fenômeno descrito na Seção \ref{sec_cl_flim}: anotar mais não implica melhorar.

\subsection{Limite de uso dos números da aplicação}

A qualidade exibida na tela é calculada contra a anotação de referência do conjunto de dados, que existe porque se trata de ambiente de demonstração. Em uso real não haveria essa referência, e o especialista julgaria pela própria máscara.

A aplicação traz ainda um guia de anotação, reproduzido na Figura~\ref{fig_plat_tutorial}, que mostra lado a lado três formas de gastar a mesma quantidade de tinta na mesma imagem. Ele existe porque a orientação de onde tocar é contraintuitiva: a tendência natural de quem anota é preencher o centro do objeto, e o guia mostra o traço encostado no contorno.

\begin{figure}[htb]
\centering
\includegraphics[width=\linewidth]{plataforma_tutorial.jpg}
\caption[Guia de anotação da plataforma]{O guia de anotação da aplicação. Mesma imagem e quantidade de pixels comparável nas três colunas; o que muda é a posição do traço. \textbf{Os valores de $F_\beta$ mostrados são ilustrativos e não constituem evidência}: vêm de uma execução única sobre três imagens, sem repetição por semente e sem teste pareado. A evidência sobre posição de anotação está no Capítulo~\ref{cap_4}.}
\label{fig_plat_tutorial}
\end{figure}

\textbf{Nenhum número produzido pela aplicação é usado nesta dissertação.} Ela roda com orçamento pequeno e execução única, regime em que a amplitude do braço de sorteio, entre 0,10 e 0,22 de $F_\beta$, é maior do que qualquer efeito que se queira observar. A aplicação serve para tornar o fenômeno visível e para permitir que um especialista opere o laço, não para medir. Todas as medições citadas vêm das campanhas em lote registradas no Capítulo~\ref{cap_4}.

Os valores apresentados pela aplicação não são usados em nenhuma afirmação desta dissertação. Eles vêm de execução única, sem repetição e sem controle de partição, e servem para inspeção qualitativa. Todas as medições relatadas no Capítulo \ref{cap_4} vêm das campanhas em lote, com semente, partição e registro.

\section{Protocolo experimental}

\subsection{Partições}

Cada conjunto de dados é dividido em três partes disjuntas: um conjunto de onde saem as imagens a anotar, um conjunto de validação usado nas medições de diagnóstico, e um conjunto de teste. A divisão é feita uma vez, com semente fixa, e mantida em todas as campanhas, de modo que resultados de campanhas diferentes sejam comparáveis.

O conjunto de teste não participou de nenhuma decisão de método: nenhum hiperparâmetro, critério, limiar ou faixa do filtro de área foi escolhido com base nele. Convém, porém, ser preciso sobre o que isso significa e o que não significa.

As campanhas de diagnóstico, que são as que sustentam os resultados centrais deste trabalho, medem \textbf{exclusivamente em validação}. É o caso do ganho marginal, da caracterização do comportamento dos critérios, da troca do gerador de candidatos e da análise de custo do aprendizado ativo.

A campanha de seleção de imagens, por outro lado, reporta desempenho \textbf{no conjunto de teste} para todos os braços, o que corresponde a cinco critérios, quatro orçamentos e três conjuntos de dados. Todos esses braços foram declarados antes da execução, de modo que se trata de reporte de configurações pré-especificadas e não de busca sobre o teste. Ainda assim, dizer que o conjunto de teste foi \emph{tocado uma única vez} seria impreciso, e essa formulação não é usada aqui.

\subsection{Repetições e unidade de análise}

Cada condição é repetida com dez sementes. Uma semente determina quais imagens são sorteadas para anotação e quais traços são gerados sobre elas, mantendo fixos a partição e o conjunto de validação.

A unidade independente de análise é a semente. Quando a mesma condição é avaliada com mais de um decodificador, os decodificadores de uma mesma semente compartilham o codificador e não constituem observações independentes; eles são colapsados por média antes de qualquer teste. Tratá-los como independentes multiplicaria artificialmente o tamanho amostral.

A consequência dessa escolha merece registro explícito: os intervalos de confiança descrevem a variação sob sorteio das imagens de treino \emph{naquele} conjunto de validação, e não generalizam para o conjunto de dados como um todo.

\subsection{Testes estatísticos}

As comparações são pareadas. Dentro de cada célula, os braços compartilham a partição, as imagens e os traços, de modo que a diferença entre eles isola o tratamento. Usa-se o teste $t$ pareado bicaudal, com nível de significância de $0{,}05$.

Quando uma campanha produz múltiplas comparações, aplica-se a correção de Benjamini-Hochberg sobre a família de testes, e o número de comparações que sobrevivem é reportado junto com os valores não corrigidos.

Resultados sobre proporções, como a taxa de sucesso do laço interativo, usam o teste exato de McNemar sobre os pares discordantes, já que pares concordantes não carregam informação sobre a diferença entre os braços.

\subsection{Pré-especificação versionada}
\label{sec_prereg}

Cada campanha é precedida de um documento que declara, antes da execução: a hipótese, o desfecho primário, os desfechos secundários, o critério que falsearia a hipótese, o número de repetições, e o que o experimento não autorizará afirmar.

O documento é versionado no mesmo repositório do código, e a ordem entre a declaração e a execução fica registrada nos \emph{commits}. Trata-se, portanto, de \textbf{pré-especificação versionada internamente}, e não de registro em plataforma externa e imutável como as usadas em ensaios clínicos ou em psicologia experimental. A distinção importa: um repositório sob controle do próprio autor oferece garantia mais fraca que um registro externo, ainda que o histórico torne a reescrita detectável. O que o procedimento garante é que a hipótese e o critério de falseamento existiam, em arquivo datado, antes de o resultado ser conhecido.

Essa disciplina tem motivo prático. Em um espaço com muitos critérios, conjuntos de dados e orçamentos, encontrar uma combinação aparentemente favorável é quase certo mesmo quando nenhum efeito existe. Declarar antes o que contaria como evidência é o que distingue um achado de um acidente. Quatro resultados aparentemente favoráveis foram descartados por esse procedimento ao longo deste trabalho, conforme relatado no Capítulo \ref{cap_4}.

Análises não previstas no pré-registro são permitidas e reportadas, mas identificadas como exploratórias, e qualquer afirmação derivada delas exige confirmação em repetições inéditas.

\subsection{Procedência dos resultados}

Toda execução avaliada é gravada em um registro único, com uma linha por execução, contendo a configuração, a semente, as imagens usadas, as métricas obtidas e o identificador do commit que produziu o resultado. O identificador de cada execução é derivado dos campos que a definem, de modo que reexecutar a mesma configuração produz o mesmo identificador e a duplicata se torna detectável.

As tabelas desta dissertação são geradas a partir desse registro por um programa, sem etapa manual. Cada tabela gerada acompanha um arquivo que lista os identificadores das execuções que produziram cada linha, o que torna a pergunta sobre a origem de um número respondível de forma direta.

\subsection{Reprodutibilidade numérica}

A estimação dos filtros envolve agrupamento por $k$-médias, cuja implementação recorre a bibliotecas de álgebra linear que distribuem o cálculo em várias linhas de execução. A ordem em que as somas parciais são reduzidas varia entre execuções, o que altera o resultado no último dígito representável.

No regime em que o agrupamento não ocorre, descrito pela Equação \ref{eq_regimes}, essa variação não tem efeito. No regime em que ocorre, ela é suficiente para que o agrupamento convirja para uma partição diferente, e o efeito se amplifica entre camadas, já que cada camada recebe como entrada as ativações da anterior.

Todas as campanhas são executadas com as bibliotecas limitadas a uma única linha de execução, condição sob a qual duas execuções idênticas produzem resultados idênticos. A verificação desse comportamento é relatada no Capítulo \ref{cap_4}.

\section{Materiais}

\subsection{Conjuntos de dados}

Três conjuntos são usados, escolhidos por apresentarem propriedades distintas quanto ao tamanho relativo do objeto, que a investigação mostrou ser determinante.

O conjunto de \textbf{parasitas} contém imagens de microscopia de ovos de \textit{Schistosoma mansoni}, em cores, com resolução de $400 \times 400$ pixels. O objeto ocupa cerca de 3\% da imagem. Estão disponíveis marcadores desenhados por dois especialistas, usados nas comparações que envolvem anotação real.

O conjunto \textbf{BraTS} \cite{menze2014multimodal} contém cortes de ressonância magnética cerebral, em tons de cinza, com resolução de $240 \times 240$ pixels.

O conjunto de \textbf{conjuntivite} contém fotografias oculares em cores e alta resolução, nas quais o objeto ocupa fração bem maior da imagem, com área mediana de aproximadamente 89 mil pixels.

A Figura~\ref{fig_datasets} mostra um exemplo de cada um, com a anotação de referência sobreposta e a máscara isolada. Três propriedades ficam visíveis de imediato e sustentam leituras feitas adiante: os três são de segmentação \emph{binária}, e não multiclasse; o objeto ocupa fração pequena da imagem em dois deles, o que é a razão de a acurácia passar de 0,97 em quase todas as configurações sem distinguir braço nenhum; e os domínios são visualmente muito distintos entre si, o que é precisamente o que torna informativo um método falhar nos três.

\begin{figure}[htb]
\centering
\includegraphics[width=\linewidth]{datasets.png}
\caption[Os três conjuntos de dados usados nos experimentos]{Um exemplo de cada conjunto, com a anotação de referência em verde na linha superior e a máscara isolada na linha inferior. A fração de pixels de objeto indicada em cada painel é medida na máscara do exemplo mostrado.}
\label{fig_datasets}
\end{figure}

\subsection{Levantamento de domínios candidatos}

A escolha dos três conjuntos acima partiu de um levantamento mais amplo de domínios em que a segmentação interativa reduziria esforço de anotação de forma significativa. O critério de seleção considerou a diversidade de desafios que cada domínio impõe, e não apenas a disponibilidade dos dados.

Em gastroenterologia e dermatologia, o Kvasir-SEG \cite{jha2020kvasir} apresenta bordas irregulares de pólipos e artefatos de iluminação, enquanto o ISIC 2018 \cite{codella2018skin} caracteriza-se por alta variabilidade de cor e textura da pele, com interferência de sombras, presença de pelos e heterogeneidade entre pacientes.

Em radiologia e oncologia, o BraTS \cite{menze2014multimodal} e o LiTS \cite{bilic2019liver} introduzem modalidades volumétricas processadas em fatias bidimensionais, com canais monocromáticos e contraste baixo entre tecidos.

Em ultrassonografia e oftalmologia, o BUSI \cite{al2020dataset} impõe ruído característico da modalidade, e o REFUGE2 \cite{fang2022refuge2} exige definição anatômica precisa de estruturas oculares.

A Figura~\ref{fig_dominios} reúne exemplos desses domínios. A diversidade de modalidade é evidente, e é justamente ela que torna o levantamento inicial plausível: endoscopia, dermatoscopia, ressonância, tomografia e ultrassonografia impõem desafios de aparência muito diferentes.

\begin{figure}[htb]
\centering
\includegraphics[width=\linewidth]{dominios_levantados.png}
\caption[Exemplos dos domínios levantados]{Exemplos dos domínios considerados no levantamento, com a anotação de referência sobreposta. A fração de objeto indicada é medida na máscara do exemplo mostrado e não é estatística do conjunto. O REFUGE2, citado no texto, não aparece porque o par de imagem e máscara disponível no repositório não corresponde entre si.}
\label{fig_dominios}
\end{figure}

Dos domínios levantados, três foram efetivamente usados nos experimentos. A razão é que a investigação revelou, ao longo do trabalho, que a propriedade determinante para o comportamento dos métodos avaliados não é a modalidade de aquisição, e sim a \emph{fração da imagem ocupada pelo objeto}, que controla a razão entre candidatos a filtro e canais disponíveis descrita na Equação \ref{eq_regimes}. Os três conjuntos escolhidos cobrem essa propriedade em faixas bem separadas, de cerca de 3\% a mais de 20\%, o que se mostrou mais informativo do que acrescentar domínios com fração semelhante.

\subsection{Modelos de comparação}

Além do \ac{FLIM} com decodificador adaptativo, foram treinados cinco modelos de referência para contextualizar os valores obtidos: SAMNet, MSCNet, MEANet, UNet e uma variante UNet com codificador \ac{FLIM}. Esses modelos são treinados por retropropagação e com conjuntos anotados completos, de modo que não constituem comparação direta quanto ao esforço de anotação; servem para situar a ordem de grandeza da qualidade atingível em cada domínio.

\subsection{Ambiente}

Os experimentos foram executados em ambiente nativo, com a biblioteca de agrupamento do \textit{scikit-learn}. Uma implementação alternativa de $k$-médias, mais rápida, está disponível na biblioteca original mas não foi utilizada, porque produz centróides diferentes e portanto filtros diferentes, o que tornaria os resultados incomparáveis entre si.

A avaliação não emprega o delineamento por árvores dinâmicas descrito no trabalho de referência, por indisponibilidade do executável na plataforma utilizada. Em seu lugar, usa-se limiarização de Otsu seguida de filtro de área, conforme descrito na Seção \ref{sec_decoder}. Os valores reproduzidos, portanto, não são diretamente comparáveis aos publicados, e essa limitação acompanha toda comparação com o trabalho original.
